\documentclass[book.tex]{subfiles}

\begin{document}
\chapter{Fiber Bundles}
\label{chap:fiber_bundles}
\noindent\rule{11cm}{0.4pt} \\
\textbf{Prerequisites:} \autoref{chap:manifold} \\
\textbf{Difficulty Level:} *** \\
\noindent\rule{11cm}{0.4pt}
\vspace{5mm}

A fiber bundle is a type of mathematical space that associates to each point on a given base space a "fiber" space. We have encountered a few simple examples of fiber bundles before, notably the tangent bundle $\mathcal{T}M$ over a manifold $M$ and the cotangent bundle $\mathcal{T}^*M$. Mathematically, a fiber bundle is defined by a projection map
\begin{align}
    \pi: E \to B
\end{align}
Here, $\pi$ is a surjection, $B$ is the base space, and $\pi$ is defined such that in local regions of $E$, $\pi$ behaves like the projection map $B \times F \to B$ where $F$ is the fiber space. We say that $\pi^{-1}(x)$ is the fiber over point $x$.

\section{Examples}

The formal mathematical definition of a fiber bundle is complex so we start with examples to build understanding.

\subsection{Trivial Bundle}

Let $E = B \times F$ and let $\pi$ be the projection onto the first coordinate. In this case, $\pi: E \to B$ is called the trivial bundle.

\subsection{Vector Bundle}

In the case that the fiber $F$ is a vector space, we say that the fiber bundle is a vector bundle. The tangent bundle $\mathcal{T}M$ and cotangent bundle $\mathcal{T}^*M$ are vector bundles.

\subsection{Principle Bundle}

%\textcolor{red}{TODO: Define a group action.}

Suppose we have a fiber bundle $\pi: E \to B$. Let $G$ be a group and suppose we have a continuous right group action
\begin{align}
P \times G \to P
\end{align}
such that the group action preserves the fibers. That is, if $y \in \pi^{-1}(x)$ then for all $g \in G$, we have $y \cdot g \in \pi^{-1}(x)$.

%\textcolor{red}{TODO: Do we want to define frame bundles here?}

\section{Section}

A section of a fiber bundle is a map
\begin{align}
    f: B \to E
\end{align}
such that $\pi(f(x)) = x$. That is, a section maps every point on the base space to a point in the fiber "above" it.

%\section{Exercises}
%\begin{enumerate}
%    \item TODO
%\end{enumerate}

\end{document}